%%%%%%%%%%%%%%%%%%%%%%%%%%%%%%%%%
%
%       EXERCÍCIO
%
%%%%%%%%%%%%%%%%%%%%%%%%%%%%%%%%%

\ifdefstring{\atividade}{prova}{%
    \ifdefstring{\modo}{objetivo}{%
        \renewcommand{\valorquestao}{\ValorQObj\ ponto}
    }{%
        \renewcommand{\valorquestao}{\ValorQDisc\ pontos}
    }%
}%

\begin{exercicioBanco}[\valorquestao]
Uma cisterna de 6000 litros foi esvaziada em um período de 3h. Na primeira hora foi utilizada apenas uma bomba, mas nas duas horas seguintes, a fim de reduzir o tempo de esvaziamento, outra bomba foi ligada junto com a primeira. O gráfico, formado por dois segmentos de reta, mostra o volume de água presente na cisterna, em função do tempo

%
\begin{center}
    \begin{tikzpicture}[scale=0.9]
        % Eixos
        \draw[->] (0,0) -- (3.5,0) node[below right] {\small{Tempo (horas)}};
        \draw[->] (0,0) -- (0,3.5) node[above left] {\small{Volume (litros)}};
              
        % Pontos escolhidos
        \draw[dashed] (1,0) -- (1,2.5) -- (0,2.5);
             
        % Marcando os pontos
        \filldraw (0,3) circle (1pt) node[left] {\small\(6000\)};
        \filldraw (1,2.5) circle (1pt) node[xshift=-41pt] {\small\(5000\)};
        \filldraw (1,0) circle (1pt) node[below] {\small\(1\)};
        \filldraw (3,0) circle (1pt) node[below] {\small\(3\)};
        
        \draw[thick, blue] (0,3) -- (1,2.5) -- (3,0);
    \end{tikzpicture}
\end{center}
%

%
Encontre a vazão, em litro por hora, da bomba que foi ligada no início da segunda hora.
%

% Define as alternativas
\newcommand{\alternativas}{%
    \begin{center}
        \begin{tabularx}{\textwidth}{XXXXX}
            (a) \(500\). &
            (b) \(1500\). &
            (c) \(2000\). &
            (d) \(2500\). &
            (e) \(3000\).
        \end{tabularx}
    \end{center}
}

% Resposta correta
\newcommand{\resposta}{B}

% Lógica condicional para exibição
\ifdefstring{\atividade}{lista}{%
    \alternativas
    \vspace{0.5em}
    
    \noindent\textbf{Resposta:} letra \textbf{\resposta}.
}{%
    \ifdefstring{\modo}{objetivo}{%
        \alternativas
    }{%
        % modo = discursiva → não mostra alternativas
    }
}
\end{exercicioBanco}

